\section{Προτάσεις που χρειάζονται απόδειξη}
\begin{Protashmabox}{Μονοτονία αντίστροφης}
Αν μία συνάρτηση $f : A \to \mathbb{R}$ είναι γνησίως αύξουσα (αντίστοιχα γνησίως φθίνουσα)
στο $A$, τότε η συνάρτηση $f^{-1}$ είναι γνησίως αύξουσα (αντίστοιχα γνησίως φθίνουσα)
στο $f(A)$.
\end{Protashmabox}
\apodeixi

Θα αποδείξουμε ότι η $f^{-1}$ είναι γνησίως αύξουσα στο $f(A)$ με \textbf{απαγωγή σε άτοπο}.

Έστω ότι η $f^{-1}$ \emph{δεν} είναι γνησίως αύξουσα στο $f(A)$. Τότε υπάρχουν
$y_1, y_2 \in f(A)$ με $y_1 < y_2$, τέτοια ώστε:
\[
  f^{-1}(y_1) \geq f^{-1}(y_2).
\]
Επειδή η $f$ είναι γνησίως αύξουσα και $f^{-1}(y_1) \geq f^{-1}(y_2)$,
εφαρμόζοντας την $f$ και στα δύο μέλη (η μονοτονία διατηρεί την ανισότητα) παίρνουμε:
\[
  f\!\left(f^{-1}(y_1)\right) \geq f\!\left(f^{-1}(y_2)\right).
\]
Από τον ορισμό της αντίστροφης, $f\!\left(f^{-1}(y)\right) = y$, οπότε:
\[
  y_1 \geq y_2,
\]
κάτι που είναι \textbf{άτοπο}, αφού υποθέσαμε $y_1 < y_2$. Συνεπώς η $f^{-1}$ είναι γνησίως αύξουσα στο $f(A)$.


% -------------------------------------------------------
% Proposition 6.2
% -------------------------------------------------------
\begin{Protashmabox}{Σημεία τομής των $C_f$, $C_{f^{-1}}$, με $f$ γν. αύξουσα}

Αν η συνάρτηση $f : \mathbb{R} \to \mathbb{R}$ είναι γνησίως αύξουσα, τότε:
\[
  f^{-1}(x) = f(x) \;\Leftrightarrow\; f(x) = x.
\]
\end{Protashmabox}

\apodeixi

Αποδεικνύουμε το ορθό και το αντίστροφο:

\begin{itemize}
  \item[$(\Rightarrow)$] Έστω $x_0 \in \mathbb{R}$ τέτοιο ώστε $f^{-1}(x_0) = f(x_0)$.
  Εφαρμόζουμε $f$ και στα δύο μέλη:
  \[
    f\!\left(f^{-1}(x_0)\right) = f\!\left(f(x_0)\right)
    \;\Rightarrow\; x_0 = f\!\left(f(x_0)\right). \tag{$\star$}
  \]
  Θα δείξουμε ότι αναγκαστικά $f(x_0) = x_0$, αποκλείοντας τις άλλες περιπτώσεις:
  \begin{itemize}
    \item Αν $f(x_0) < x_0$: Επειδή $f$ γνησίως αύξουσα, εφαρμόζουμε $f$:
          $f\!\left(f(x_0)\right) < f(x_0)$.
          Από $(\star)$: $x_0 < f(x_0)$, \textbf{άτοπο}.
    \item Αν $f(x_0) > x_0$: Ομοίως,
          $f\!\left(f(x_0)\right) > f(x_0)$,
          και από $(\star)$: $x_0 > f(x_0)$, \textbf{άτοπο}.
  \end{itemize}
  Άρα αναγκαστικά $f(x_0) = x_0$.

  \item[$(\Leftarrow)$] Έστω $x_0 \in \mathbb{R}$ τέτοιο ώστε $f(x_0) = x_0$.
  Επειδή $f(x_0) = x_0$, από τον ορισμό της αντίστροφης:
  \[
    f(x_0) = x_0 \;\Rightarrow\; f^{-1}(x_0) = x_0 = f(x_0),
  \]
  δηλαδή $f^{-1}(x_0) = f(x_0)$.
\end{itemize}

Άρα τελικά:
\[
  f^{-1}(x) = f(x) \;\Leftrightarrow\; f(x) = x,
\]
δηλαδή τα σημεία τομής των $C_f$ και $C_{f^{-1}}$ ταυτίζονται με τα σημεία τομής του $C_f$
με τη διχοτόμο $y = x$.


% -------------------------------------------------------
% Proposition 6.3
% -------------------------------------------------------
\begin{Protashmabox}{Όριο απόλυτης τιμής}
Ισχύει ότι:
\[
  \lim_{x \to x_0} f(x) = 0 \;\Leftrightarrow\; \lim_{x \to x_0} |f(x)| = 0.
\]
\end{Protashmabox}
\apodeixi

\begin{itemize}
  \item[$(\Rightarrow)$] Αν $\displaystyle\lim_{x \to x_0} f(x) = 0$, τότε χρησιμοποιούμε
  την ιδιότητα ορίων για αριθμητικές πράξεις (η απόλυτη τιμή είναι συνεχής συνάρτηση):
  \[
    \lim_{x \to x_0} |f(x)| = \left|\lim_{x \to x_0} f(x)\right| = |0| = 0.
  \]

  \item[$(\Leftarrow)$] Αν $\displaystyle\lim_{x \to x_0} |f(x)| = 0$, χρησιμοποιούμε
  την ανισότητα που ισχύει για κάθε $x \in D_f$:
  \[
    -|f(x)| \leq f(x) \leq |f(x)|.
  \]
  Παίρνουμε όρια και στα τρία μέλη:
  \[
    \lim_{x \to x_0}\!\left(-|f(x)|\right) = -0 = 0
    \quad \text{και} \quad
    \lim_{x \to x_0} |f(x)| = 0.
  \]
  Από το \textbf{κριτήριο παρεμβολής}, επειδή και τα δύο άκρα τείνουν στο $0$, συμπεραίνουμε:
  \[
    \lim_{x \to x_0} f(x) = 0.
  \]
\end{itemize}


% -------------------------------------------------------
% Proposition 6.4
% -------------------------------------------------------
\begin{Protashmabox}{Όριο της $f^2$}
Ισχύει ότι:
\[
  \lim_{x \to x_0} f^2(x) = 0 \;\Leftrightarrow\; \lim_{x \to x_0} f(x) = 0.
\]
\end{Protashmabox}
\apodeixi
\vspace{-5mm}
\begin{itemize}
  \item[$(\Rightarrow)$] Αν $\displaystyle\lim_{x \to x_0} f(x) = 0$, τότε από τις ιδιότητες
  ορίων (γινόμενο ορίων):
  \[
    \lim_{x \to x_0} f^2(x) = \left(\lim_{x \to x_0} f(x)\right)^2 = 0^2 = 0.
  \]

  \item[$(\Leftarrow)$] Αν $\displaystyle\lim_{x \to x_0} f^2(x) = 0$, παρατηρούμε ότι
  $\sqrt{f^2(x)} = |f(x)|$, οπότε:
  \[
    \lim_{x \to x_0} |f(x)| = \lim_{x \to x_0}\sqrt{f^2(x)} = \sqrt{0} = 0.
  \]
  Χρησιμοποιούμε και την ανισότητα που ισχύει για κάθε $x \in D_f$:
  \[
    -|f(x)| \leq f(x) \leq |f(x)|.
  \]
  Επειδή $\displaystyle\lim_{x \to x_0}(-|f(x)|) = 0$ και
  $\displaystyle\lim_{x \to x_0}|f(x)| = 0$, από το \textbf{κριτήριο παρεμβολής}:
  \[
    \lim_{x \to x_0} f(x) = 0.
  \]
\end{itemize}


% -------------------------------------------------------
% Proposition 6.5
% -------------------------------------------------------
\begin{Protashmabox}{$1-1\Leftrightarrow$ γν. μονότονη}
Αν μία συνάρτηση $f : \mathbb{R} \to \mathbb{R}$ είναι συνεχής και $1-1$, τότε η $f$ είναι γνησίως μονότονη.
\end{Protashmabox}
\apodeixi

Θα αποδείξουμε με \textbf{απαγωγή σε άτοπο}. Έστω ότι η $f$ \emph{δεν} είναι
γνησίως μονότονη. Τότε υπάρχουν τιμές:
\[
  x_1 < x_2 < x_3 \quad \text{με:} \quad f(x_1) > f(x_2) \quad \text{και} \quad f(x_2) < f(x_3).
\]
(Δηλαδή η $f$ δεν είναι ούτε γνησίως αύξουσα ούτε γνησίως φθίνουσα στο σύνολο ορισμού της.)

Επιλέγουμε έναν αριθμό $\gamma$ ανάμεσα στο $f(x_2)$ και στο ελάχιστο των $f(x_1)$, $f(x_3)$:
\[
  f(x_2) < \gamma < \min\{f(x_1),\, f(x_3)\}.
\]
Τότε:
\begin{itemize}
  \item Στο διάστημα $[x_1, x_2]$: ισχύει $f(x_2) < \gamma < f(x_1)$, και επειδή $f$ συνεχής,
        από το \textbf{Θεώρημα Ενδιάμεσων Τιμών} υπάρχει $\alpha \in (x_1, x_2)$ με $f(\alpha) = \gamma$.
  \item Στο διάστημα $[x_2, x_3]$: ισχύει $f(x_2) < \gamma < f(x_3)$, και ομοίως
        υπάρχει $\beta \in (x_2, x_3)$ με $f(\beta) = \gamma$.
\end{itemize}

Έτσι έχουμε $f(\alpha) = \gamma = f(\beta)$, ενώ $\alpha \neq \beta$
(αφού $\alpha < x_2 < \beta$). Αυτό \textbf{αντιφάσκει} με την υπόθεση ότι η $f$ είναι $1-1$.

Άρα η $f$ είναι γνησίως μονότονη.


\begin{Protashmabox}{Ολοκλήρωμα άρτιας και περιττής}
Έστω συνάρτηση $f:[-a,a]\to\mathbb{R}$, η οποία είναι συνεχής. Να δείξετε ότι:
\begin{alist}
  \item Αν η $f$ είναι περιττή, τότε:
        $\displaystyle\int_{-a}^{a} f(x)\,dx = 0$.
  \item Αν η $f$ είναι άρτια, τότε:
        $\displaystyle\int_{-a}^{a} f(x)\,dx = 2\int_{0}^{a} f(x)\,dx$.
\end{alist}
\end{Protashmabox}
\apodeixi

Και στις δύο περιπτώσεις, σπάμε το ολοκλήρωμα στα $[-a,0]$ και $[0,a]$:
\[
  \int_{-a}^{a} f(x)\,dx = \int_{-a}^{0} f(x)\,dx + \int_{0}^{a} f(x)\,dx.
\]
Το κλειδί είναι να μετασχηματίσουμε το $\displaystyle\int_{-a}^{0} f(x)\,dx$
με την αντικατάσταση $u = -x$ (οπότε $du = -dx$):
\begin{itemize}
  \item Όταν $x = -a$: $u = a$.
  \item Όταν $x = 0$: $u = 0$.
\end{itemize}
Οπότε:
\[
  \int_{-a}^{0} f(x)\,dx
  \;\overset{u=-x}{=}\;
  \int_{a}^{0} f(-u)\cdot(-du)
  = \int_{0}^{a} f(-u)\,du
  = \int_{0}^{a} f(-x)\,dx.
\]

\begin{alist}
  \item \textbf{Περιττή $f$:} Ισχύει $f(-x) = -f(x)$, οπότε
  $\displaystyle\int_{-a}^{0} f(x)\,dx = \int_{0}^{a} f(-x)\,dx = -\int_{0}^{a} f(x)\,dx$.
  Άρα:
  \[
    \int_{-a}^{a} f(x)\,dx
    = -\int_{0}^{a} f(x)\,dx + \int_{0}^{a} f(x)\,dx = 0.
  \]
  Αυτό είναι αναμενόμενο γεωμετρικά: η γραφική παράσταση περιττής συνάρτησης
  έχει συμμετρία ως προς την αρχή, οπότε τα «εμβαδά» αλληλοαναιρούνται.

  \item \textbf{Άρτια $f$:} Ισχύει $f(-x) = f(x)$, οπότε
  $\displaystyle\int_{-a}^{0} f(x)\,dx = \int_{0}^{a} f(-x)\,dx = \int_{0}^{a} f(x)\,dx$.
  Άρα:
  \[
    \int_{-a}^{a} f(x)\,dx
    = \int_{0}^{a} f(x)\,dx + \int_{0}^{a} f(x)\,dx
    = 2\int_{0}^{a} f(x)\,dx.
  \]
  Αυτό αντικατοπτρίζει τη συμμετρία ως προς τον άξονα $y$: το εμβαδό για $x < 0$
  ισούται με αυτό για $x > 0$.
\end{alist}


% -------------------------------------------------------
% Proposition 6.8
% -------------------------------------------------------
\begin{Protashmabox}{Μηδενικό ολοκλήρωμα}
Έστω $f : [a,\beta] \to \mathbb{R}$ μια συνάρτηση, η οποία είναι συνεχής και ισχύει
$f(x) \geq 0$, για κάθε $x \in [a,\beta]$. Να δείξετε ότι:
\[
  \int_{a}^{\beta} f(x)\,dx = 0 \;\Leftrightarrow\; f(x) = 0, \quad x \in [a,\beta].
\]
\end{Protashmabox}
\apodeixi
\vspace{-3mm}
\begin{itemize}
  \item[$(\Leftarrow)$] Αν $f(x) = 0$ για κάθε $x \in [a,\beta]$, τότε ολοκληρώνουμε
  τη μηδενική συνάρτηση, οπότε προφανώς:
  \[
    \int_{a}^{\beta} f(x)\,dx = \int_{a}^{\beta} 0\,dx = 0.
  \]

  \item[$(\Rightarrow)$] Έστω ότι $\displaystyle\int_{a}^{\beta} f(x)\,dx = 0$. Θα δείξουμε
  ότι αναγκαστικά $f(x) = 0$ για κάθε $x \in [a,\beta]$, με {απαγωγή σε άτοπο}.

  Έστω ότι υπάρχει $x_0 \in [a,\beta]$ ώστε $f(x_0) \neq 0$.
  Επειδή $f(x) \geq 0$, αυτό σημαίνει $f(x_0) > 0$.
  Επειδή η $f$ είναι \textbf{συνεχής} στο $x_0$, για αρκετά μικρό διάστημα
  γύρω από το $x_0$, παραμένει $f(x) > 0$.
  Ένα θετικό ολοκλήρωμα σε αυτό το υποδιάστημα «συνεισφέρει» θετικά, και
  αφού $f(x) \geq 0$ παντού, δεν μπορεί να αντισταθμιστεί. Άρα:
  \[
    \int_{a}^{\beta} f(x)\,dx > 0,
  \]
  κάτι που έρχεται σε αντίθεση με την υπόθεση $\displaystyle\int_{a}^{\beta} f(x)\,dx = 0$. Άρα, $f(x) = 0$ για κάθε $x \in [a,\beta]$.
\end{itemize}
